\documentclass{article}
\usepackage{hyperref}
\usepackage{amsmath}
\usepackage{amssymb}
\usepackage{todonotes} % disable this package in final thesis

\title{MINLP Modeling}
\author{Bjarne Jehn}


\newcommand{\noam}{noam} %numberOfAppearingMinimizers

\newcommand{\timeForCheckInR}{{T_R}}
\newcommand{\timeForLookupInS}{{T_{SL}}}
\newcommand{\timeForHashtable}{{T_H}}

\newcommand{\percKmers}{\textbf{p}_\textbf{K}}
\newcommand{\percKmersEmpiric}{P_{K}}

\newcommand{\possibleMinimizerLength}{\mathcal{M}}
\newcommand{\possibleThresholds}{\mathcal{T}}

\newcommand{\minimizerLength}{\textbf{m}}
\newcommand{\threshold}{\textbf{t}}

\newcommand{\selectorm}{\textbf{sm}}
\newcommand{\selectort}{\textbf{st}}

\newcommand{\levelI}{{L_i}}
\newcommand{\levelOne}{{L_1}}
\newcommand{\levelTwo}{{L_2}}
\newcommand{\levelThree}{{L_3}}

\begin{document}
	
\maketitle
\newpage



\section*{One Level}
\subsection*{Runtime}
As a first approach we model the runtime complexity for only one level.\\
In the following we will use capital letters for constants and bold letters for variables.\\
\\
\begin{tabular}{ l | l }
	constants	& meaning \\
	\hline
	$\timeForCheckInR$		& time to check presence of a bit in vector R \\
	$\timeForLookupInS$		& time to lookup a k-mer in vector S on level i \\
	$\timeForHashtable$		& time to lookup a k-mer on level i \\
	$\percKmersEmpiric(m, t)$ &empirically estimated value for a kmer with a minimizer \\
	& of length m to have a frequency less than t \\
	$\possibleMinimizerLength$			& Set of all minimizer lengths to be considered \\
	$\possibleThresholds$			& Set of thresholds to be considered \\
	\hline
	\hline
	variables	&  \\
	\hline
	$\percKmers \in \left[0,1\right]$ & probability of a kmer to have a minimizer with frequency $\leq t$ \\
	$\selectorm_m \in \left\{0, 1\right\}$	& binary selector for minimizer length m\\
	$\selectort_t \in \left\{0, 1\right\}$	& binary selector for threshold of t\\
\end{tabular}\\
\\
For one level the lookup of a query with minimizer m of frequency f(m) would consist of checking the bitvector R in order to determin whether $1 \leq f(m) \leq threshold$. This will be the case for $\percKmers$ of the k-mers, which will then be followed by a lookup in S. If m does not appear at all or has a $f(m) > threshold$, a lookup in the hashtable will follow. The lookup time for one level can thus be modeled as:
\begin{equation*}
	\timeForCheckInR + \timeForLookupInS^\levelOne * \percKmers + \timeForHashtable * (1-\percKmers)
\end{equation*}
\\
Let $\possibleMinimizerLength$ be the Set of all minimizer length and $\possibleThresholds$ the set of all thresholds to be considered.\\
In order to determine $m$ and $t$ we introduce binary selector variables $\selectorm_m$, $\selectort_t$ for each possible $m \in \possibleMinimizerLength$ and $t \in \possibleThresholds$.\\
$\selectorm_m = 
	\begin{cases}
		1 & \text{minimizer length is } m\\
		0 & \text{otherwise}.
	\end{cases}$\\
$\selectort_t = 
\begin{cases}
	1 & \text{threshold is } t\\
	0 & \text{otherwise}.
\end{cases}$\\
For a fixed minimizer length $m$ and threshold $t$ let $\percKmersEmpiric(m, t)$ be the empirically determined percentage of k-mers that have a minimizer with $frequency < t$.\\
The selectors for m and t must be binary and mutually exclusive:
\begin{align*}
	&\selectorm_m \in {0,1} ~|~ \forall m \in \possibleMinimizerLength \\
	&\selectort_t \in {0,1} ~|~ \forall t \in \possibleThresholds \\
	&\sum_{m \in \possibleMinimizerLength} \selectorm_m = 1 \\
	&\sum_{t \in \possibleThresholds} \selectort_t = 1
\end{align*}
We can then model $\percKmers$ with:
\begin{align*}
	&\percKmers = \sum_{m \in \possibleMinimizerLength} \left( \selectorm_m \sum_{t \in \possibleThresholds} \percKmersEmpiric(m, t) * \selectort_t\right)\\
	\implies &\percKmers - \sum_{m \in \possibleMinimizerLength} \left( \selectorm_m \sum_{t \in \possibleThresholds} \percKmersEmpiric(m, t) * \selectort_t\right)= 0
\end{align*}
$\minimizerLength$ and $\threshold$ can then be determined by:
\begin{align*}
	\minimizerLength &= \sum_{m \in \possibleMinimizerLength} \selectorm_m * m \\
	\threshold &= \sum_{t \in \possibleThresholds} \selectort_t * t
\end{align*}
\\
The above results in the following MINLP (Mixed Integer Non Linear Program):\\
\\
MINIMIZE
\begin{align*}
	&\timeForCheckInR + \timeForLookupInS^\levelOne * \percKmers + \timeForHashtable * (1-\percKmers)
\end{align*}
SUBJECT TO
\begin{align*}
		&\percKmers - \sum_{m \in M} \left( \selectorm_m \sum_{t \in T} \percKmersEmpiric(m, t) * \selectort_t\right) = 0\\
		&\selectorm_m \in {0,1} ~|~ \forall m \in M \\
		&\selectort_t \in {0,1} ~|~ \forall t \in T \\
		&\sum_{m \in M} \selectorm_m = 1 \\
		&\sum_{t \in T} \selectort_t = 1 \\
		&\minimizerLength = \sum_{m \in \possibleMinimizerLength} \selectorm_m * m \\
		&\threshold = \sum_{t \in \possibleThresholds} \selectort_t * t
\end{align*}
\\
For now we will assume that $\timeForLookupInS \approx 20ns$ and $\timeForHashtable \approx 40ns$.\\
Solving this MINLP will result in a Solution where the $\selectorm_m$, $\selectort_t$ corresponding to the largest $m \in \possibleMinimizerLength$, $t \in \possibleThresholds$ will be set to one.\\
This is as expected, since this leads to the highest possible $\percKmers$. This maximizes the k-mers found in the vector $S$ and thus minimizes the lookups in the (slower) Hashtable.\\
\\
However this also results in a high space consumption, because \todo{explain why}\\
\\
In order to take that into account, we have to add a constraint modeling the space consumption.\\

\subsection*{Space Consumption}
The space consumption of the Data Structure used by RSHash can be split into that of the Vectors $R$, $O$ and $S$ as well as the hashtable $H$.

\subsubsection*{Bitvectors $R$ and $S$}
$R$ and $S$ can likely be disregarded when modeling the space consumption. \todo{explain why}

\subsubsection*{Offsetvector $O$}
$O$ contains one entry of size $log |T|$ for the offset of every minimizer with frequency $\leq t$ in $T$.\\
\[s_O = log_2|T| * |O| = log_2|T| * \sum_{m=1}^{t} \minimizerLength * \text{(number of minimizers with frequenzy m)} \]
Every minimizer of length $\minimizerLength$ covers an average of $\frac{k-\minimizerLength+2}{2}$ consecutive k-mers. We can thus model the bits per k-mer in $O$ as:\\
\[bpk_O = \frac{2}{k-\minimizerLength+2} * \percKmers * log |T| \]

\subsubsection*{Hashtable $H$}
$H$ contains one entry of size $log |T|$ for the offset of every minimizer with frequency $> t$ in $T$.\\
However $H$ also contains a number of empty buckets. \todo{reason}\\
For a hashtable with $b$ buckets containing $n$ elements, the load factor $\alpha$ is defined as:
\[\alpha = \frac{n}{b}\].\\
That means that for every entry in the hashtable, there must be $\frac{1-\alpha}{\alpha}$ many empty buckets.\\
Using this we can model the bits per k-mer in $H$ as:\\
\[bpk_H = (\alpha + \frac{1-\alpha}{\alpha}) * \frac{2}{k-\minimizerLength+2} * (1-\percKmers) * log |T| \]


\newpage
\section*{Several Levels}
\subsection*{Runtime}
minimize: nicht linear!\\


\[
\begin{array}{llll}
	\timeForCheckInR^\levelOne 	&+ \percKmers^\levelOne*\timeForLookupInS^\levelOne 				&							&\\
			&+ (1-\percKmers^\levelOne) * \Bigg( \timeForCheckInR^\levelTwo &+ \percKmers^\levelTwo*\timeForLookupInS^\levelTwo 				&\\
			&							&+ (1-\percKmers^\levelTwo) * \Big( \timeForCheckInR^\levelThree	&+ \percKmers^\levelThree*\timeForLookupInS^\levelThree \\
			&							&							&+ (1-\percKmers^\levelThree) * \timeForHashtable \Big) \Bigg)
\end{array}
\]
\\
TODO: linearise
\\
\\
\begin{tabular}{ l | l }
	variable	& meaning \\
	\hline
	$\percKmers^\levelI$		& probability of a kmer belonging to level i \\
	$\timeForCheckInR^\levelI$		& time to check presence of a bit in vector R of level i \\
	$t_l^\levelI$		& time to lookup a k-mer on level i \\
\end{tabular}
\\
\\

We can model $\percKmers$ for multiple levels with:
\begin{align*}
	&\percKmers^\levelI = \sum_{m \in M} \left( \selectorm_m^\levelI \sum_{t \in T} \percKmersEmpiric^\levelI(m, t) * \selectort_t^\levelI\right)\\
	&\percKmers^\levelI - \sum_{m \in M} \left( \selectorm_m^\levelI \sum_{t \in T} \percKmersEmpiric^\levelI(m, t) * \selectort_t^\levelI\right)= 0
\end{align*}
\\
\\
TODO: linearise
\\
\\

\begin{tabular}{ l | l }
	variable	& meaning \\
	\hline
	$\selectorm_m^\levelI$	& binary selector for minimizer length m in level i\\
	$\selectort_t^\levelI$	& binary selector for threshold of t in level i\\
	$\percKmersEmpiric^\levelI(m, t)$ &empirically estimated value for a kmer with a minimizer \\
				& of length m to have a frequency less than t \\
	$M$			& Set of all minimizer length to be considered \\
	$T$			& Set of thresholds to be considered \\
\end{tabular}
\\
\\
selectors for m and t must be binary amd mutually exclusive:
\begin{align*}
	\selectorm_m^\levelI &\in {0,1} ~|~ \forall m \in M \\
	\selectort_t^\levelI &\in {0,1} ~|~ \forall t \in T \\
	\sum_{m \in M} \selectorm_m^\levelI &= 1 \\
	\sum_{t \in T} \selectort_t^\levelI &= 1 \\
\end{align*}

Without further constraints

\subsection*{Space Consumption} \label{sec:space}

Using Elias Fano Coding a Bitvector of length $U$ with $n$ bits set to 1 can be encoded using $n \lceil log \frac{U}{n} \rceil + 3n$.\\
\\
Let $\noam$ be the number of appearing minimizers (minimizers with a frequency greater than zero).

\subsubsection*{Level 1}
$R$ - Presence of k-mers:\\
\textbf{probably not relevant}
$R^\levelOne$ has length $4^{sm_m^\levelOne}$ and contains $p_M^\levelOne * \noam$ many 1s.\\

\begin{align*}
	s_R^\levelOne 
	&= p_M^\levelOne * \noam \lceil log \frac{4^{sm_m^\levelOne}}{p_M^\levelOne * \noam} \rceil + 3 * p_M^\levelOne *  \noam\\
	&= p_M^\levelOne * \noam \left( \lceil log \frac{4^{sm_m^\levelOne}}{p_M^\levelOne * \noam} \rceil + 3 \right)
\end{align*}
\\
\\
$S$ - Offset Boundaries:\\
$S$ is a bitvector of length $|O| + 1$ and is \textbf{probably not relevant} for the optimization of $m$ and $t$.
\\
\\
$O$ - Offsets:\\
$O^\levelOne$ contains one entry of size $log |T|$ for the offset of every minimizer with frequency $^\leq t^\levelOne$ in $T$.\\
\[s_O = log_2|T| * |O| = log_2|T| * \sum_{m=1}^{t^\levelOne} m * \text{(number of minimizers with frequenzy m)} \]
Every minimizer of length $m$ covers an average of $\frac{k-m+2}{2}$ consecutive k-mers.\\
\[bpk_O^\levelOne = \frac{2}{k-m^\levelOne+2} * p_K^\levelOne * log |T| \]


\end{document}